\EMhold

\EMsection{Approximating a Periodic Function by a Finite Trigonometric Series}{Approximating a Periodic Function by a Finite Trigonometric Series}

\begin{EMtheorem}{}

For an \EMim{N}-term approximation of a periodic function \EMim{f(x)} by a series of sines and cosines, the best coefficients (in the MMSE sense) are the Fourier series coefficients.

\end{EMtheorem}

Consider a periodic function with period \EMim{2\pi}:
\EMdisp{f(x)\EMbr =a_0+\sum_{n=1}^{\infty}a_n\cos nx+b_n\sin nx,\qquad\EMbr  f(x)\EMbr =f(x+2\pi)}
The \EMim{N}-term trigonometric approximation:
\EMdisp{f(x)\EMbr \approx a_0+\sum_{n=1}^{N}a_n\cos nx+b_n\sin nx}
Is the following trigonometric approximation not \textbf{better}?
\EMdisp{F(x)\triangleq\alpha_0+\sum_{n=1}^{N}\alpha_n\cos nx+\beta_n\sin nx}
What is the \textbf{criterion} for an approximation being better? One criterion: the error function:
\EMdisp{E(x)\EMbr =\big|f(x)-F(x)\big|\EMbr =\left|f(x)-\alpha_0-\sum_{n=1}^{N}\big(\alpha_n\cos nx+\beta_n\sin nx\big)\right|}
A good criterion: the mean squared error:
\EMdisp{D\triangleq\int_{-\pi}^{\pi}E^2(x)\,dx\EMbr =\int_{-\pi}^{\pi}\left|f(x)-\alpha_0-\sum_{n=1}^{N}\big(\alpha_n\cos nx+\beta_n\sin nx\big)\right|^2dx}

\begin{EMproof}{}

For the best approximation we minimize the mean squared error. This is done by differentiating \EMim{D} with respect to the coefficients of the trigonometric series and setting these derivatives equal to zero. In other words:
\EMdisp{\frac{\partial D}{\partial\alpha_0}\EMbr =-2\int_{-\pi}^{\pi}\left[f(x)-\alpha_0-\sum_{n=1}^{N}\big(\alpha_n\cos nx+\beta_n\sin nx\big)\right]dx\EMbr =0}
\EMdisp{\int_{-\pi}^{\pi}f(x)\,dx\EMbr =\int_{-\pi}^{\pi}\left[\alpha_0+\sum_{n=1}^{N}\big(\alpha_n\cos nx+\beta_n\sin nx\big)\right]dx\EMbr =2\pi\alpha_0\quad\EMbr \EMbr \Longrightarrow\quad\EMbr \alpha_0\EMbr =\frac1{2\pi}\int_{-\pi}^{\pi}f(x)\,dx\EMbr =a_0}
\EMdisp{\frac{\partial D}{\partial\alpha_m}\EMbr =0\quad\EMbr \EMbr \Longrightarrow\quad\EMbr \alpha_m\EMbr =\frac1\pi\int_{-\pi}^{\pi}f(x)\cos mx\,dx\EMbr =a_m}
\EMdisp{\frac{\partial D}{\partial\beta_m}\EMbr =0\quad\EMbr \EMbr \Longrightarrow\quad\EMbr \beta_m\EMbr =\frac1\pi\int_{-\pi}^{\pi}f(x)\sin mx\,dx\EMbr =b_m}
The best finite trigonometric approximation of a periodic function is the finite Fourier series.

\end{EMproof}

\EMhold

\EMsection{Differentiation and Integration of Fourier Series}{Differentiation and Integration of Fourier Series}

\begin{EMtheorem}{Differentiating a Fourier series}

If \EMim{f(x)} is periodic with period \EMim{T}, continuous and has a Fourier series, then its derivative is also periodic, and its Fourier series is obtained by differentiating the Fourier series of the function.
\EMdisp{f(x)\EMbr =a_0+\sum_{n=1}^{\infty}a_n\cos\frac{2\pi}Tnx+b_n\sin\frac{2\pi}Tnx}
\EMdisp{\Longrightarrow\quad\EMbr  f'(x)\EMbr =\sum_{n=1}^{\infty}\left(\frac{2\pi}Tnb_n\right)\cos\frac{2\pi}Tnx+\left(-\frac{2\pi}Tna_n\right)\sin\frac{2\pi}Tnx}

\end{EMtheorem}

\begin{EMtheorem}{Integrating a Fourier series}

If \EMim{f(x)} is periodic with period \EMim{T} and has a Fourier series, then the definite integral of this function can be computed by integrating the Fourier series of the function.
\EMdisp{f(x)\EMbr =a_0+\sum_{n=1}^{\infty}a_n\cos\frac{2\pi}Tnx+b_n\sin\frac{2\pi}Tnx}
\EMdisp{\Longrightarrow\quad\EMbr \int_0^xf(t)\,dt\EMbr =a_0x+\sum_{n=1}^{\infty}\left(\frac T{2\pi n}a_n\sin\frac{2\pi}Tnx-\frac T{2\pi n}b_n\cos\frac{2\pi}Tnx+\frac T{2\pi n}b_n\right)}

\end{EMtheorem}

\EMhold

\EMsection{Rayleigh's Theorem}{Rayleigh’s Theorem}

\begin{EMtheorem}{Rayleigh's theorem}

If the functions \EMim{f(x)} and \EMim{g(x)} are periodic with period \EMim{T} and their Fourier series exist as
\EMdisp{f(x)\EMbr =a_0+\sum_{n=1}^{\infty}a_n\cos\frac{2\pi}Tnx+b_n\sin\frac{2\pi}Tnx}
\EMdisp{g(x)\EMbr =\alpha_0+\sum_{n=1}^{\infty}\alpha_n\cos\frac{2\pi}Tnx+\beta_n\sin\frac{2\pi}Tnx}
then
\EMdisp{\frac2T\int_Tf(x)\,\overline{g(x)}\,dx\EMbr =2a_0\overline{\alpha_0}+\sum_{n=1}^{\infty}a_n\overline{\alpha_n}+b_n\overline{\beta_n}}

\end{EMtheorem}

\begin{EMproof}{}

Proof, assuming \EMim{g(x)} is real:
\EMdisp{\begin{aligned}
\frac2T\int_Tf(x)g(x)\,dx&=\frac2T\int_Tf(x)\left(\alpha_0+\sum_{n=1}^{\infty}\alpha_n\cos\frac{2\pi}Tnx+\beta_n\sin\frac{2\pi}Tnx\right)dx\\
&=\frac2T\int_Tf(x)\alpha_0\,dx+\sum_{n=1}^{\infty}\frac2T\int_Tf(x)\left(\alpha_n\cos\frac{2\pi}Tnx+\beta_n\sin\frac{2\pi}Tnx\right)dx\\
&=\underbrace{\frac2T\int_Tf(x)\,dx}_{2a_0}\,\alpha_0+\sum_{n=1}^{\infty}\underbrace{\frac2T\int_Tf(x)\cos\frac{2\pi}Tnx\,dx}_{a_n}\,\alpha_n+\underbrace{\frac2T\int_Tf(x)\sin\frac{2\pi}Tnx\,dx}_{b_n}\,\beta_n\\
&=2a_0\alpha_0+\sum_{n=1}^{\infty}a_n\alpha_n+b_n\beta_n
\end{aligned}}

\end{EMproof}

\EMhold

\EMsection{Parseval's Theorem}{Parseval’s Theorem}

\begin{EMtheorem}{Parseval's theorem}

If the function \EMim{f(x)} is periodic with period \EMim{T} and its Fourier series exists as
\EMdisp{f(x)\EMbr =a_0+\sum_{n=1}^{\infty}a_n\cos\frac{2\pi}Tnx+b_n\sin\frac{2\pi}Tnx}
then
\EMdisp{\frac2T\int_T|f(x)|^2\,dx\EMbr =2|a_0|^2+\sum_{n=1}^{\infty}\big(|a_n|^2+|b_n|^2\big)}

\end{EMtheorem}

\begin{EMproof}{}

In Rayleigh's identity, replacing \EMim{g(x)} by \EMim{f(x)} gives Parseval's identity.

\end{EMproof}

By Parseval's identity, the power of the function \EMim{f(x)} can be obtained from the Fourier coefficients. If the function is expressed by an \EMim{N}-term approximation of sines and cosines, then by the relation above the power of the approximating function is less than the power of \EMim{f(x)}, and we have:

\begin{EMimportant}{Bessel's inequality}

\EMdisp{\frac2T\int_T|f(x)|^2\,dx\ \ge\ 2|a_0|^2+\sum_{n=1}^{N}\big(|a_n|^2+|b_n|^2\big)}

\end{EMimportant}

\begin{EMexample}{}

For the function \EMim{f(x)} below, how many terms of the Fourier series are needed for at least 94 percent of the power to remain in the approximating function?

\EMdisp{f(x)\EMbr =\sum_{n=1}^{\infty}\frac{2k}{n\pi}(1-\cos n\pi)\sin nx}
\EMdisp{\frac2T\int_T|f(x)|^2\,dx\EMbr =\frac1\pi\int_{-\pi}^{\pi}k^2\,dx\EMbr =\frac2\pi\big[k^2x\big]_0^\pi\EMbr =2k^2}
\EMdisp{\frac{\displaystyle\sum_{n=1}^{N}\left(\frac{2k}{n\pi}(1-\cos n\pi)\right)^2}{2k^2}\ge0.94}
\EMdisp{\frac{\displaystyle\sum_{n=1}^{7}\left(\frac{2k}{n\pi}(1-\cos n\pi)\right)^2}{2k^2}\EMbr \cong0.9496,\qquad\EMbr  \frac{\displaystyle\sum_{n=1}^{5}\left(\frac{2k}{n\pi}(1-\cos n\pi)\right)^2}{2k^2}\EMbr \cong0.9331\quad\EMbr \EMbr \Longrightarrow\quad\EMbr  N\EMbr =7}

\end{EMexample}

\begin{EMunit}

\EMfigure{m02-power-example}{The square wave \EMim{\pm k} with period \EMim{2\pi} used in the power
example}

\end{EMunit}

\EMsection{The Complex Form of the Fourier Series}{The Complex Form of the Fourier Series}

Suppose \EMim{f(x)} is periodic and, merely for simplicity, has period \EMim{2\pi}. Then
\EMdisp{f(x)\EMbr =a_0+\sum_{n=1}^{\infty}a_n\cos nx+b_n\sin nx}
Using Euler's formula \EMim{e^{ix}=\cos x+i\sin x}:
\EMdisp{\cos nx\EMbr =\frac{e^{inx}+e^{-inx}}2,\qquad\EMbr  \sin nx\EMbr =\frac{e^{inx}-e^{-inx}}{2i},\qquad\EMbr  \frac1i\EMbr =-i}
\EMdisp{f(x)\EMbr =a_0+\sum_{n=1}^{\infty}a_n\left(\frac{e^{inx}+e^{-inx}}2\right)+b_n\left(\frac{e^{inx}-e^{-inx}}{2i}\right)\EMbr =\underbrace{a_0}_{c_0}+\sum_{n=1}^{\infty}\underbrace{\left(\frac{a_n-ib_n}2\right)}_{c_n}e^{inx}+\underbrace{\left(\frac{a_n+ib_n}2\right)}_{c_{-n}}e^{-inx}}
\EMdisp{c_n\EMbr =\frac{a_n-ib_n}2\EMbr =\frac12\left\{\frac1\pi\int_{2\pi}f(x)\cos nx\,dx-i\frac1\pi\int_{2\pi}f(x)\sin nx\,dx\right\}\EMbr =\frac1{2\pi}\int_{2\pi}f(x)\big(\cos nx-i\sin nx\big)\,dx\EMbr =\frac1{2\pi}\int_{2\pi}f(x)e^{-inx}\,dx}
\EMdisp{c_{-n}\EMbr =\frac{a_n+ib_n}2\EMbr =\frac1{2\pi}\int_{2\pi}f(x)\big(\cos nx+i\sin nx\big)\,dx\EMbr =\frac1{2\pi}\int_{2\pi}f(x)e^{inx}\,dx}

\begin{EMimportant}{The complex form of the Fourier series}

\EMdisp{f(x)\EMbr =\sum_{n=-\infty}^{\infty}c_ne^{inx},\qquad\EMbr  c_n\EMbr =\frac1{2\pi}\int_{2\pi}f(x)e^{-inx}\,dx}
For an arbitrary period:
\EMdisp{f(x)\EMbr =\sum_{n=-\infty}^{\infty}c_ne^{i\frac{2\pi}Tnx},\qquad\EMbr  c_n\EMbr =\frac1T\int_Tf(x)e^{-i\frac{2\pi}Tnx}\,dx}
\EMdisp{\frac1T\int_T|f(x)|^2\,dx\EMbr =\sum_{n=-\infty}^{\infty}|c_n|^2}

\end{EMimportant}
